\documentclass[11pt]{article}
\usepackage[OT1]{fontenc}
\renewcommand{\rmdefault}{cmr}
\renewcommand{\sfdefault}{cmss}
\renewcommand{\ttdefault}{cmtt}
\usepackage{amsmath,amssymb,amsthm}
\usepackage[margin=1in]{geometry}
\usepackage[hidelinks]{hyperref}
\newtheorem{theorem}{Theorem}[section]
\newtheorem{lemma}[theorem]{Lemma}
\newtheorem{proposition}[theorem]{Proposition}
\newtheorem{definition}[theorem]{Definition}
\newtheorem{remark}[theorem]{Remark}
\newcommand{\N}{\mathbb N}
\newcommand{\lv}{\operatorname{lv}}
\newcommand{\norm}{\operatorname{N}}
\newcommand{\sz}{\operatorname{s}}
\newcommand{\Add}{\operatorname{Add}}
\newcommand{\EP}{\operatorname{ExpProd}}
\newcommand{\CP}{\operatorname{Psi}}
\newcommand{\Rep}{\operatorname{Repeat}}
\newcommand{\one}{\mathbf1}
\newcommand{\pa}{\partial}
\title{Literal Sum Conventions in the Wilken--Weiermann Assignment}
\author{}
\date{30 October 2026}
\begin{document}
\maketitle
\begin{abstract}
The Wilken--Weiermann abstraction operator tests whether an entire ordinal
expression is free of its bound variable before recursing into its binary
syntax. Consequently an unspecified association of finite sums cannot be
recovered from their ordinal values. We give explicit finite parsing policies,
prove the actual-source forward translation for each policy, and exhibit a
closed level-one source whose positive assigned coordinate is respectively
204 and 184 under two permitted associations. We also correct a strict
inequality in the source's scalar-partial lemma. A positive-grade subterm
estimate proves bounded-norm preservation with the original multiplier,
without using the false strict inequality. These results make a
policy-relative comparison question precise; they do not assert a failure of
termination or identify a uniquely intended historical parsing policy.
\end{abstract}

\section{Why ordinal values do not specify the input tree}

An ordinal expression in \cite[Definition~2.1]{WW} is a literal tree. Its
three binary constructors, with immediate subterms $f,g$, are
\[
 \Add(f,g),\qquad \EP(f,g),\qquad \CP(f,g).
\]
Their interpretations are $f\oplus g$, $2^f\otimes g$ and
$\psi(\omega\otimes f\oplus g)$, respectively. Here $\oplus,\otimes$
are natural sum and product. The last two expressions are primitive
constructors, rather than arbitrary multiplication trees. In this note,
ordinary arithmetic on natural-number sizes is distinguished from these
literal ordinal trees.

For a tracked variable vector $X$, the size $\sz_X(u)$ is one when the
whole $u$ is $X$-free, or when $u$ is a component of $X$. Otherwise it is
the sum of the child sizes plus one at Add and Psi, and twice the head
size plus the tail size plus one at ExpProd. The partial operator also
uses the whole-free case first. Thus
\[
 \sz_X(x_0+(1+1))=3,\qquad
 \sz_X((x_0+1)+1)=5.
\]
The ordinal values of the two sums agree. This small example alone does
not show that the distinction occurs in a generated source assignment;
Section~\ref{sec:witness} supplies such an assignment.

The source specifies finite sums and natural-number multiples in
Definitions~2.10 and2.18 without a full binary association convention.
Its convention for left-associated lambda-term application concerns a
different grammar. Also, actual graded coordinate trees are zero-free,
whereas out-of-range vector components are written as zero. Footnote5
to Definition~2.10 expressly explains that the box definition is designed
to avoid zero occurrences in coordinate parse trees. A literal
implementation must respect both facts.

\section{Explicit finite parsing policies}

\begin{definition}
An admissible policy $\rho$ is a computable rule choosing an ordered full
binary addition tree for each nonempty finite-sum macro site and its
arity. Each operand is used once as an unchanged subtree. The choice is
independent of operand names, values, internal syntax and free-variable
status. A $k$-fold repeat, $k>0$, has exactly $k$ unchanged copies of
its operand. Existing subtrees are never flattened, normalized or
reassociated.
\end{definition}

The macro sites are finite numerals; positive-integer repeats; the
three-operand active-sum partial; the three operands
$\Rep_2(\pa_{i+1}f),\pa_i g,\one$ in an active exponential partial;
the positive-coordinate abstraction sums; its scalar repeat; and the
optional three-operand box tail. Repeat2 remains a single nested
operand in the exponential clause. Binary clauses stay binary. The
integer
\[
 S_X(h)=2^n\sum_{i=0}^m\sz_X(h_i),
 \qquad n=\lv(X)+1,
\]
is computed in natural-number arithmetic from the already fixed input
trees, then used as a repeat count.

Each vector records its actual level. An out-of-range lookup has an
\emph{absent} role, interpreted as zero for comparison. A vector sum
omits absent operands but retains every actual operand unchanged. The
box's optional finite-level contribution is absent when its argument
level is zero; it is a positive numeral otherwise. This is not a rule
that simplifies actual Zero nodes or tests whether an expression has
value zero. Every actual source coordinate remains zero-free. In
particular the source program $0$ is assigned One.

A useful default takes left combs for sum, repeat, numeral and positive
aggregation macros, including the active-sum partial, while retaining
the already explicit box tail $f_0+(g_0+n)$ for $n>0$ and $f_0+g_0$
for $n=0$. This is a declared convention, not a claim that the source
uniquely selected it.

\begin{proposition}\label{prop:finite}
For a fixed policy, all literal partials and zero-history assignments
are uniquely determined finite trees. Their expansions have computable
bounds in the input trees and type levels. For a fixed input tree,
changing only the output associations of a partial does not change
its ordinal value.
\end{proposition}
\begin{proof}
Structural recursion determines a partial. Each zero-history lambda
then uses its body's fixed tree, its fixed sizes, and the finite
repeats specified by the policy. Erasing only the newly supplied
parentheses recovers the printed formulas exactly.

For an explicit bound, let $|u|$ count tree nodes. Give a leaf weight
one, Add/Psi weight $w(f)+w(g)+1$, and ExpProd weight
$2w(f)+w(g)+1$. Induction gives
\[
 \sz_X(u)\le w(u)\le2^{|u|}-1.
\]
Every coordinate of a partial has at most $E(|u|)=6\cdot3^{|u|}$ nodes.
At a whole-free diagonal it has $|u|+2$ nodes. At active Add, ExpProd,
scalar Psi and positive Psi outputs the respective bounds are
\[
 E(f)+E(g)+3,\quad 2E(f)+E(g)+4,\quad
 |f|+E(g)+1,\quad E(f)+E(g)+1,
\]
where $E(f)$ abbreviates $E(|f|)$. These fit $E(1+|f|+|g|)$.
Thus $S_X(h)\le2^n\sum_i(2^{|h_i|}-1)$. The scalar repeat has
$S_X(h)(|(\pa_0h_0)_0|+1)-1$ nodes; a positive abstraction sum has
at most $\sum_i E(|h_i|)+m$ nodes. The box is a finite downward
recursion, with $|b_i|=|b_{i+1}|+|f_i|+|g_i|+2$ at $1\le i\le n$.
These bounds do not evaluate $\psi$.

For fixed input trees, structural induction on the partial clauses
uses associativity of natural sum and the same finite multiplicities
to prove output-value independence. An enclosing partial need not see
identical input trees, so this does not prove source-policy independence.
\end{proof}

\section{The weak scalar lemma and the graded domain}

The domain $C_i^X$ is as in \cite[Definition~2.9]{WW}. Actual coordinate
trees are positive. Its scalar class has finite values; the natural norm
of a finite value is that value. Write $P_i=\pa_i^X$ and $n=\lv(X)+1$.
All expression inequalities below hold under every admissible substitution.

\begin{remark}\label{rem:error}
Printed Lemma~2.19, page~22 of \cite{WW}, asserts
$h\{X:=\one\}<(P_0h)_0$ for all $h\in C_0^X$.
At $h=x_0$ both sides equal one. This is a false strict assertion in
the actual PDF, not an extraction artifact. The appendix's strict
estimate (A.5) also fails: for $\lv(X)=0$, $h=(x_0,1)$, $i=1$ and
the free subterm $t=1$, it asserts $1<1$.
\end{remark}

\begin{theorem}[Weak scalar replacement]\label{thm:weak}
For $h\in C_0^X$,
\[
 h\{X:=\one\}\le(P_0h)_0.
\]
Every partial $P_i h$, $h\in C_i^X$, is a valid $X$-free $C$-vector
of level $n$. For a vector $h\in C^X$, the assembled $\delta^X h$ is likewise valid and $X$-free.
These assertions hold for every admissible policy.
\end{theorem}
\begin{proof}
Induct simultaneously on the grammar and the partial recursion. A
whole-free scalar has partial $h+1$, a tracked scalar variable has
partial one, and an active sum has scalar partial $P(f)+P(g)+1$.
At an active collapse the two heads after $X:=\one$ are identical,
and the induction hypothesis compares the finite tails. For a fixed
ordinal $a$ and finite $r$, the published inequality
$\psi(a)+r\le\psi(a+r)$, Proposition~2.3(1), gives the needed weak
finite-tail monotonicity. No global monotonicity of $\psi$ is used.

For the new collapse's domain premise, the original premise and the
weak induction hypothesis give
\[
 \norm(f\{X:=\one\})
 \le F_2(g\{X:=\one\})\le F_2((P_0g)_0).
\]
Its head has the correct positive grade. Whole-free diagonal outputs
are $h+1$, off-diagonal outputs and tracked-variable outputs are One.
The other cases only sum correctly graded partial outputs, with the
stated finite multiplicities. Positive collapse outputs add the head
partial and the positive tail partial. Each output is $X$-free.
Assembling the finite sums and the positive scalar repeat preserves
the domain. A coordinate above $n$ is already $X$-free: there is no
tracked component at its grade, and positive promotion of a scalar
requires whole-$X$-freeness. It may be retained unchanged.
\end{proof}

\section{Bounded norm with the original multiplier}

The additive norm satisfies the published inequalities
\[
 \norm(a\oplus b)=\norm(a)+\norm(b),\quad
 \norm(a)+\norm(b)-1\le\norm(a\otimes b)\quad(a,b>0),
 \qquad \norm(a)\le2\norm(2^a).
\]
Also $\norm(a)\le\psi(a)$ and $r\le\psi(a+r)$ for finite $r$.
These are Propositions~2.2 and2.3 of \cite{WW}.

\begin{lemma}\label{lem:free}
Let $h\in C^X$ be bounded in norm. Fix an admissible substitution for its other
names and put $B=h_0\{X:=\one\}$. Every first wholly $X$-free
subtree encountered at a positive grade in a partial recursion starting
at any actual coordinate of $h$ satisfies
\[
 \norm(t)\le2^{n-1}B.
\]
\end{lemma}
\begin{proof}
Compute norms after $X:=\one$; the free $t$ is unchanged and every
input coordinate then has norm at most $B$. Descending through an Add
or an exponential tail does not increase norm. An exponential-head
step costs at most two, since
\[
 \norm(f)\le2\norm(2^f)\le2\norm(2^f\otimes g),
\]
and increases grade by one. At grade $n$ a subtree is already wholly
$X$-free and recursion stops. From a positive root there are at most
$n-1$ such steps.

A scalar path can enter a positive grade only through a collapse head.
This costs no factor two: positivity and the cited bounds give
\[
 \norm(f)\le\norm(\omega\otimes f\oplus g)
                  \le\psi(\omega\otimes f\oplus g),\qquad
 g\le\psi(\omega\otimes f\oplus g).
\]
The head enters grade one, leaving at most $n-1$ active
exponential-head steps. Scalar tails stay at grade zero. A promoted
scalar at a positive grade is already whole-free, so cannot cause an
additional active descent. This proves the bound.
\end{proof}

\begin{theorem}\label{thm:norm}
For $h\in C^X$ of bounded norm, $\delta^X h$ has bounded norm with the original
multiplier $S_X(h)=2^n\sum_i\sz_X(h_i)$, for every admissible policy.
\end{theorem}
\begin{proof}
Continue at an arbitrary fixed substitution for the other names. Set
$H=(P_0h_0)_0$, so $1\le B\le H$ by Theorem~\ref{thm:weak}, and
put $\nu=2^nH$. A free positive-grade subtree which can be copied into
a positive diagonal output satisfies
\[
 \norm(t)+1\le2^{n-1}B+1\le2^nH=\nu.
\]
A free scalar subtree is not copied into a positive output; those
outputs are One. This distinction pays the extra diagonal unit without
the source's false strict estimate (A.5).

Induct inside each partial to obtain, for $0<j\le n$,
\[
 \norm((P_i h_i)_j)\le\sz_X(h_i)\nu.
\]
The whole-free and tracked-variable cases follow from the preceding
paragraph. For Add the sum of the two bounds plus one is at most
$(\sz_X(f)+\sz_X(g)+1)\nu$. For ExpProd there are two copies of the
head bound, one tail bound and One, giving
$(2\sz_X(f)+\sz_X(g)+1)\nu$. A positive collapse output only adds
the head and tail bounds. Norm additivity makes all associations valid.
Summing over $i$ gives
\[
 \norm((\delta^Xh)_j)\le\Bigl(\sum_i\sz_X(h_i)\Bigr)2^nH
                      =(\delta^Xh)_0.
\]
For $j>n$ the unchanged $h_j$ is $X$-free, so its norm is at most
$B\le H\le(\delta^Xh)_0$. This holds at every substitution.
\end{proof}

\section{An actual closed source separates the policies}\label{sec:witness}

Let L and R agree at every macro except the active-sum partial, where
they choose respectively $(a+b)+1$ and $a+(b+1)$. Repeats in both
are left combs. Consider the literal source
\[
 A=\lambda z:o.\lambda x:o.((\lambda w:o.z)\,x)
                  :o\to o\to o.
\]
It is closed and recursor-free; every internal type has level at most
one. There is no beta normalization before assignment.

The assignment of $\lambda w.z$ is $(f(z),1)$, with the literal tree
\[
 f(z)=((z+1)+(z+1))+z,\qquad\sz_z(f)=9.
\]
The x-body scalar is $g=\CP(1,f(z)+x)$, whose x-size is five and
whose abstraction multiplier is ten. Put
\[
 u=f+1,\quad T_L=(u+1)+1,\quad T_R=u+(1+1),\quad
 K_\rho=\CP(1,T_\rho),\quad B_*=\CP(1,u).
\]
The x-abstraction has
\[
 h_0=\Rep_{10}(K_\rho)+B_*,\qquad
 h_1=(1+1)+\operatorname{Sum3}_\rho(1,1,1).
\]
The latter is wholly z-free and has value five. Literal sizes and
positive partial values are as follows:
\[
\begin{array}{c|rrrrrrr}
 &f&u&T_\rho&K_\rho&\Rep_{10}(K_\rho)&B_*&h_0\\\hline
 \sz_z\text{ under L}&9&11&15&17&179&13&193\\
 \sz_z\text{ under R}&9&11&13&15&159&13&173
\end{array}
\]
At every scalar expression in this table, output one of its z-partial
has the corresponding size as its finite value. Whole-free scalar
inputs give one, a tracked variable gives one, active sums add the
two values and one, and each collapse adds the head partial's value
two to its tail's value. Thus the positive output of $P_0h_0$ is
193 or 173. Since $(P_1h_1)_1$ has value six, the outer abstraction gives
\[
 \boxed{\ [\![A]\!]_{1,L}=193+6+5=204,
 \qquad [\![A]\!]_{1,R}=173+6+5=184.\ }
\]
These are distinct assigned ordinal values, not merely distinct trees.

The scalar values also separate. Write $H(t)=\psi(\omega+t)$. Direct
scalar partial calculation gives $P_0f=11$, $P_0u=14$, and
$P_0T_L=20$, $P_0T_R=18$. The outer size multipliers are 388 and 348;
the common retained value is $10H(8)+H(6)$. Consequently
\begin{align*}
 [\![A]\!]_{0,L}&=388(10H(20)+H(14)+10)+10H(8)+H(6),\\
 [\![A]\!]_{0,R}&=348(10H(18)+H(14)+10)+10H(8)+H(6).
\end{align*}
The source's finite-tail inequality gives $H(20)>H(18)$, so the L
scalar is strictly larger too. No numerical evaluation of the control
or collapse is needed. The accompanying bounded AST control reconstructs
the typed source and confirms the finite 204/184 coordinates exactly.

\section{What the forward translation preserves}

\begin{theorem}\label{thm:translation}
Fix an admissible policy. For a literal original source all of whose free
term names are ground, every positive final coordinate is closed and
effectively evaluable. Its scalar coordinate translates effectively,
preserving every shared valuation, to an expression in positive-natural
inputs, positive constants, addition, and fixed functions
\[
 h_\lambda(t)=\psi(\lambda+t),\qquad
 \lambda=\omega\otimes c>0\quad(c\text{ closed}).
\]
Arbitrary internal binders, types and recursors are allowed.
\end{theorem}
\begin{proof}
Induct simultaneously on source assignments and partial recursions.
A generated positive field contains only positive-index components of
variables, positive constants, Add and the graded ExpProd constructor.
In particular it contains no component of index zero of any name.
Primes and positive box fields preserve this stronger invariant.
A positive whole-free partial copies only its positive input at the
diagonal and returns constants elsewhere. Active-variable outputs are
constant; Add and ExpProd partials preserve the invariant. A scalar
partial's positive outputs are constants in the whole-free and variable
cases, recurse at sums, and at a collapse add a positive-head partial
and a positive output of its scalar-tail partial. These are exactly the
simultaneous induction hypotheses. Abstraction sums these outputs and
removes all components of its bound name, including in its retained part.

The only positive-index variables that could remain are positive
components of free source names. There are none when those names are
ground. All final positive fields, including scalar-collapse heads,
are therefore closed. Their positivity follows from the graded domain
proved above. Scalar generated fields contain only ground components,
constants, Add and Psi. Replace each closed-head Psi by the displayed
$h_\lambda$ applied to its recursively translated tail. The syntactic
size multipliers are fixed integers, expanded as finite additions; they
do not introduce fresh independent inputs. This proves the exact
shared-valuation statement.
\end{proof}

Only the recursor prime introduces $\omega$ into generated positive
fields. Thus in a recursor-free source the final positive fields are
finite. If all internal types and subtypes have level at most two,
coordinate two is built by sums from $\omega$ and finite constants,
so is $a\omega+b$. Coordinate one then lies in $\N[\omega]$, since
$2^{a\omega+b}=\omega^a2^b$. Every direct scalar index is consequently
below $\omega^\omega$. These assertions hold for each policy; the
example above shows that they do not imply policy independence.

\section{Limits of the repair and later uses of strictness}

Theorem~\ref{thm:weak} pays the domain use immediately following printed
Lemma~2.19. Theorem~\ref{thm:norm} replaces the false strict (A.5) in
Appendix~A.2 and proves its required (A.6) and bounded-norm conclusion
with the same constant. Its higher-coordinate use needs only
$\norm(h_j)\le B\le H\le(\delta h)_0$.

Appendix~A.7(2) also invokes the false strict link
$F_2(\bar g)<F_2((P_0g)_0)$. In its active-collapse induction, let
$Q=\psi(\omega\sz_X(g)\alpha+\sz_X(g)m)$ be the strict upper bound
on $(P_0g)_0$ from the induction hypothesis. The repaired chain is
\[
 \norm(\bar f)\le F_2(\bar g)
       \le F_2((P_0g)_0)<F_2(Q).
\]
The needed strict slack is supplied by that upper-bound induction
hypothesis, not by the false lemma. This describes that local repair;
it is not a fresh audit of every quantitative appendix estimate.

The distinct beta supplier, Lemma~2.22, compares
$h<(P_i h\square X)_i$. Its variable case is strict because of the
box and the extra input. It does not assert the false strict scalar
lemma without the box. The original termination proof uses
history-dependent assignment derivations. It does not assert that
fixed zero-history values decrease under all reductions.

The results here establish exact, policy-relative domains, norms and
the forward translation. They do not supply a full comparison decision
procedure, reconstruct a unique historical policy, or assert that all
extended annotated-recursion histories have the stronger generated
invariant of Theorem~\ref{thm:translation}.

\begin{thebibliography}{1}
\bibitem{WW}
G.~Wilken and A.~Weiermann,
\emph{Derivation lengths classification of G\"odel's~T extending
Howard's assignment},
Logical Methods in Computer Science \textbf{8}(1:19) (2012), 1--44.
\href{https://doi.org/10.2168/LMCS-8(1:19)2012}{doi:10.2168/LMCS-8(1:19)2012}.
\href{https://arxiv.org/abs/1203.0115}{arXiv:1203.0115}.
\end{thebibliography}
\end{document}
